\begin{solution}{125}
    \begin{subsolution}
        中子以初速度与静止的靶核发生碰撞，碰撞前瞬间初速度方向与两球球心连线之间的夹角为$\alpha$。在两球心连线和中子初速度方向所决定的平面上，令$x$轴沿两球心连线，中子初速度垂直于连线的方向为$y$轴。设碰撞后中子的速率为$v$，沿着$x$轴方向的速度分量为$v _ { x }$，$y$轴方向的速度分量为$v _ { y }$，碰撞后靶核速率为$v _ { 1 }$，碰撞前后沿$x$轴和$y$轴方向的动量分别守恒
        \begin{equation}
            \left\{ \begin{array}{l} m v _ {0} \cos \alpha = m v _ {x} + m _ {1} v _ {1}, \\ m v _ {0} \sin \alpha = m v _ {y}. \end{array} \right.
            \label{eq:1.s125}
        \end{equation}
        式中
        \begin{equation}
            v _ {x} ^ {2} + v _ {y} ^ {2} = v ^ {2}
            \label{eq:2.s125}
        \end{equation}
        \bigskip
        \noindent [解法（二）

        由碰撞前后动量守恒，碰前中子的动量$m v _ { 0 }$、碰后中子的动量$m v$和碰后原子核的动量$m _ { 1 } v _ { 1 }$构成一闭合的矢量三角形，如解题图1a所示。据余弦定理有
        \begin{equation}
            m ^ {2} v _ {0} ^ {2} + m _ {1} ^ {2} v _ {1} ^ {2} - 2 m m _ {1} v _ {0} v _ {1} \cos \alpha = m ^ {2} v ^ {2}
            \label{eq:3.s125}
        \end{equation}
        式中，$\alpha$是碰前中子的动量与碰后原子核的动量之间的夹角。]

        能量守恒给出
        \begin{equation}
            \frac {1}{2} m v _ {0} ^ {2} = \frac {1}{2} m v ^ {2} + \frac {1}{2} m _ {1} v _ {1} ^ {2}
            \label{eq:4.s125}
        \end{equation}
        由\eqref{eq:1.s125}和\eqref{eq:4.s125}式得
        \begin{equation}
            m _ {1} v _ {1} ^ {2} + \frac {m _ {1} ^ {2}}{m} v _ {1} ^ {2} - 2 m _ {1} v _ {0} v _ {1} \cos \alpha = 0,
            \label{eq:5.s125}
        \end{equation}
        由此得
        \begin{equation}
            v _ {1} = \frac {2 m \cos \alpha}{m + m _ {1}} v _ {0},
            \label{eq:6.s125}
        \end{equation}
        将\eqref{eq:6.s125}式代入\eqref{eq:4.s125}式得
        \begin{equation}
            v = \frac {\sqrt {(m + m _ {1}) ^ {2} - 4 m m _ {1} \cos^ {2} \alpha}}{m + m _ {1}} v _ {0} = \frac {\sqrt {m ^ {2} + m _ {1} ^ {2} - 2 m m _ {1} \cos 2 \alpha}}{m + m _ {1}} v _ {0}
            \label{eq:7.s125}
        \end{equation}
        由\eqref{eq:6.s125}式知，当$\alpha = 0$时$v _ { 1 }$达到最大，
        \begin{equation}
            v _ {1 \mathrm{max}} = \frac {2 m}{m + m _ {1}} v _ {0},
            \label{eq:8.s125}
        \end{equation}
        所以氢核的最大速率是
        \begin{equation}
            v _ {\mathrm{H}} = \frac {2 m}{m + m _ {\mathrm{H}}} v _ {0},
            \label{eq:9.s125}
        \end{equation}
        氮核的最大速率是
        \begin{equation}
            v _ {\mathrm{N}} = \frac {2 m}{m + 14 m _ {\mathrm{H}}} v _ {0},
            \label{eq:10.s125}
        \end{equation}
        由以上两式得
        \begin{align}
            m &= \frac {14 v _ {\mathrm{N}} - v _ {\mathrm{H}}}{v _ {\mathrm{H}} - v _ {\mathrm{N}}} m _ {\mathrm{H}} = \frac {14 \times 4.7 \times 10^{6} - 3.3 \times 10^{7}}{3.3 \times 10^{7} - 4.7 \times 10^{6}} m _ {\mathrm{H}} = 1.16 m _ {\mathrm{H}}, \label{eq:11.s125} \\
            v _ {0} &= \frac {m + m _ {\mathrm{H}}}{2 m} v _ {\mathrm{H}} = 3.07 \times 10^{7} \mathrm{m/s}. \label{eq:12.s125}
        \end{align}
    \end{subsolution}
    \begin{subsolution}
        速度为$V _ { i }$的氮14核继续与速度为$v _ { 0 }$的第$i$个中子碰撞，在每次碰撞后获得最大速率增量条件下，氮14核的速度变为$V _ { i + 1 }$，中子的末速度为$v _ { 0 } ^ { \prime }$，由动量守恒和能量守恒有
        \begin{equation}
            \left\{ \begin{array}{l} m v _ {0} + m _ {\mathrm{N}} V _ {i} = m v _ {0} ^ {\prime} + m _ {\mathrm{N}} V _ {i + 1}, \\ \frac {1}{2} m v _ {0} ^ {2} + \frac {1}{2} m _ {\mathrm{N}} V _ {i} ^ {2} = \frac {1}{2} m v _ {0} ^ {\prime 2} + \frac {1}{2} m _ {\mathrm{N}} V _ {i + 1} ^ {2} \end{array} \right.
            \label{eq:13.s125}
        \end{equation}
        由\eqref{eq:13.s125}式得
        \begin{equation}
            \frac {V _ {i}}{v _ {0}} = 1 - a + a \frac {V _ {i - 1}}{v _ {0}}
            \label{eq:14.s125}
        \end{equation}
        式中
        \begin{equation}
            a = \frac {m _ {\mathrm{N}} - m}{m _ {\mathrm{N}} + m} = 0.847
            \label{eq:15.s125}
        \end{equation}
        按\eqref{eq:14.s125}式逐次迭代得
        \begin{equation}
            \begin{array}{r l} \frac {V _ {n}}{v _ {0}} & = 1 - a + a \frac {V _ {n - 1}}{v _ {0}} = 1 - a + a (1 - a + a \frac {V _ {n - 2}}{v _ {0}}) = 1 - a ^ {2} + a ^ {2} \frac {V _ {n - 2}}{v _ {0}} \\ & = 1 - a ^ {3} + a ^ {3} \frac {V _ {n - 3}}{v _ {0}} = \dots = 1 - a ^ {n} + a ^ {n} \frac {V _ {0}}{v _ {0}} = 1 - a ^ {n} \end{array}
            \label{eq:16.s125}
        \end{equation}
        这里，应用了题给条件$V _ {0} = 0$。所求的次数$n$满足
        \begin{equation}
            \frac {1}{2} m _ {\mathrm{N}} V _ {n} ^ {2} = \frac {1}{2} \times 1.16 m _ {\mathrm{H}} v _ {0} ^ {2}
            \label{eq:17.s125}
        \end{equation}
        即是
        \begin{equation}
            \frac {1}{2} 14 m _ {\mathrm{H}} (1 - a ^ {n}) ^ {2} v _ {0} ^ {2} = \frac {1}{2} \times 1.16 m _ {\mathrm{H}} v _ {0} ^ {2},
            \label{eq:18.s125}
        \end{equation}
        由\eqref{eq:15.s125}和\eqref{eq:18.s125}式得，满足方程的最接近的值是
        \begin{equation}
            n = 2.
            \label{eq:19.s125}
        \end{equation}
    \end{subsolution}
    \begin{subsolution}
        根据麦克斯韦速率分布
        \begin{equation}
            f (v) = 4 \pi \left(\frac {m}{2 \pi k T}\right) ^ {3 / 2} \mathrm{e} ^ {- \frac {m v ^ {2}}{2 k _ {\mathrm{B}} T}} v ^ {2}
            \label{eq:20.s125}
        \end{equation}
        速率取极大值的条件是
        \begin{equation}
            \frac {\mathrm{d} f (v)}{\mathrm{d} v} = 0,
            \label{eq:21.s125}
        \end{equation}
        可知最概然速率为
        \begin{equation}
            v _ {\mathrm{p}} = \sqrt {\frac {2 k _ {\mathrm{B}} T}{m}}
            \label{eq:22.s125}
        \end{equation}
        最概然速率对应的动能为
        \begin{equation}
            E _ {\mathrm{p}} = \frac {1}{2} m v _ {\mathrm{p}} ^ {2} = k _ {\mathrm{B}} T.
            \label{eq:23.s125}
        \end{equation}
        设总粒子数为$N$，由动能分布函数定义可知为
        \begin{equation}
            f \left(E _ {\mathrm{k}}\right) = \frac {\mathrm{d} N}{N \mathrm{d} E _ {\mathrm{k}}}
            \label{eq:24.s125}
        \end{equation}
        由速率分布函数的定义可知
        \begin{equation}
            f (v) = \frac {\mathrm{d} N}{N \mathrm{d} v}
            \label{eq:25.s125}
        \end{equation}
        而
        \begin{equation}
            \mathrm{d} v = \frac {\mathrm{d} E _ {\mathrm{k}}}{\sqrt {2 m E _ {\mathrm{k}}}}
            \label{eq:26.s125}
        \end{equation}
        所以
        \begin{equation}
            f \left(E _ {\mathrm{k}}\right) = \frac {\mathrm{d} N}{N \mathrm{d} E _ {\mathrm{k}}} = \frac {1}{\sqrt {2 m E _ {\mathrm{k}}}} f \left(\sqrt {\frac {2 E _ {\mathrm{k}}}{m}}\right)
            \label{eq:27.s125}
        \end{equation}
        联立上述各式得
        \begin{equation}
            f (E _ {\mathrm{k}}) = \frac {2 \pi}{(\pi k _ {\mathrm{B}} T) ^ {3 / 2}} \mathrm{e} ^ {- \frac {E _ {\mathrm{k}}}{k _ {\mathrm{B}} T}} E _ {\mathrm{k}} ^ {1 / 2}
            \label{eq:28.s125}
        \end{equation}
        动能取极大值的条件为$\frac { \mathrm { d } f ( \mathcal { E } _ { \mathrm { k } } ) } { \mathrm { d } \mathcal { E } _ { \mathrm { k } } } = 0$，由此可知最概然动能为
        \begin{equation}
            E _ {\mathrm{kp}} = \frac {1}{2} k _ {\mathrm{B}} T = \frac {1}{2} E _ {\mathrm{p}}.
            \label{eq:29.s125}
        \end{equation}
    \end{subsolution}
\end{solution}